\section{Term-structure models}

\subsection{Introduction}

Fixed-income markets quote bond prices rather than a single directly observable interest rate. A coupon-bearing bond can be decomposed into zero-coupon cash flows, and the prices of those zero-coupon bonds determine yields and forward rates.

Let $P(t,T)$ denote the time-$t$ price of a zero-coupon bond that pays one unit of currency at maturity $T$, where $0\leq t\leq T$. Its continuously compounded yield to maturity is
\[
y(t,T)
=-\frac{\log P(t,T)}{T-t},
\qquad t<T,
\]
so that
\[
P(t,T)=e^{-y(t,T)(T-t)}.
\]

If a bond pays deterministic cash flows $C_i$ at dates
\[
t<T_1<\cdots<T_n,
\]
then its no-arbitrage price is
\[
\Pi_t=\sum_{i=1}^n C_iP(t,T_i).
\]
Observed coupon-bond prices can therefore be used to bootstrap zero-coupon prices. For example, if bond $j$ pays $C_{j,i}$ at $T_i$, with $C_{j,j}\neq0$, then the triangular system
\[
\Pi_j=\sum_{i=1}^j C_{j,i}P(t,T_i),
\qquad j=1,\ldots,n,
\]
can be solved successively for $P(t,T_1),\ldots,P(t,T_n)$.

The map
\[
T\longmapsto y(t,T)
\]
is the yield curve observed at time $t$. It may be upward sloping, downward sloping, flat, or humped; none of these shapes is universal.

For $t\leq T_1<T_2$, the simple forward rate fixed at time $t$ for lending over $[T_1,T_2]$ is
\[
L(t;T_1,T_2)
=\frac{1}{T_2-T_1}
 \left(\frac{P(t,T_1)}{P(t,T_2)}-1\right).
\]
Indeed, buying one $T_1$-bond and financing it by shorting
$P(t,T_1)/P(t,T_2)$ units of the $T_2$-bond creates a loan from $T_1$ to $T_2$.

The instantaneous forward rate is
\[
f(t,T)
=-\frac{\partial}{\partial T}\log P(t,T),
\]
which implies
\[
P(t,T)
=\exp\!\left(-\int_t^Tf(t,u)\,du\right).
\]
The short rate is
\[
r_t=f(t,t).
\]

There are two principal modeling approaches:
\begin{enumerate}
    \item specify the dynamics of the short rate or of a finite collection of state variables and derive bond prices; or
    \item specify the dynamics of the entire forward-rate curve, as in the Heath--Jarrow--Morton framework.
\end{enumerate}

\subsection{Affine term-structure models}

\subsubsection*{One-factor models}

The Vasicek model specifies, under a risk-neutral measure $\widetilde{\mathbb{P}}$,
\[
dr_t=\kappa(\theta-r_t)dt+\sigma\,d\widetilde W_t,
\qquad
\kappa,\theta,\sigma>0.
\]
This Gaussian model is mean reverting but permits negative rates. The time-inhomogeneous extension
\[
dr_t=\bigl(\theta(t)-\kappa r_t\bigr)dt+\sigma\,d\widetilde W_t
\]
is commonly called the one-factor Hull--White model.

The Cox--Ingersoll--Ross model is
\[
dr_t=\kappa(\theta-r_t)dt+\sigma\sqrt{r_t}\,d\widetilde W_t,
\qquad
\kappa,\theta,\sigma>0.
\]
Starting from $r_0\geq0$, the process remains nonnegative. The Feller condition
\[
2\kappa\theta\geq\sigma^2
\]
makes the boundary at zero unattainable when $r_0>0$, and hence ensures strict positivity.

Suppose generally that, under $\widetilde{\mathbb{P}}$,
\[
dr_t=a(t,r_t)dt+b(t,r_t)d\widetilde W_t.
\]
The zero-coupon bond price is
\[
P(t,T)
=\widetilde{\mathbb{E}}\!\left[
 \exp\!\left(-\int_t^Tr_sds\right)
 \bigg|\mathcal{F}_t
\right]
=v(t,r_t;T).
\]
By Feynman--Kac, $v$ solves
\[
\begin{cases}
 v_t+a(t,r)v_r+\dfrac12b(t,r)^2v_{rr}-rv=0,
    & t<T,\\[4pt]
 v(T,r;T)=1.
\end{cases}
\]
In affine models, the solution has exponential-affine form
\[
P(t,T)=A(t,T)e^{-C(t,T)r_t},
\]
so the yield
\[
y(t,T)
=-\frac{\log A(t,T)}{T-t}
 +\frac{C(t,T)}{T-t}r_t
\]
is affine in the current short rate.

\subsubsection*{Two-factor Gaussian models}

A two-factor Vasicek model may be written as
\[
\begin{aligned}
dX_1(t)
&=\bigl(a_1-b_{11}X_1(t)-b_{12}X_2(t)\bigr)dt
 +\sigma_1d\widetilde W_1(t),\\
dX_2(t)
&=\bigl(a_2-b_{21}X_1(t)-b_{22}X_2(t)\bigr)dt
 +\sigma_2d\widetilde W_2(t),\\
r_t&=\delta_0+\delta_1X_1(t)+\delta_2X_2(t),
\end{aligned}
\]
where
\[
d\widetilde W_1(t)d\widetilde W_2(t)=\rho\,dt,
\qquad -1<\rho<1.
\]
After a linear transformation, one may often use a canonical form driven by independent Brownian motions, such as
\[
\begin{aligned}
dY_1(t)&=-\lambda_1Y_1(t)dt+d\widetilde B_1(t),\\
dY_2(t)&=-\lambda_{21}Y_1(t)dt-\lambda_2Y_2(t)dt+d\widetilde B_2(t),\\
r_t&=\delta_0+\delta_1Y_1(t)+\delta_2Y_2(t).
\end{aligned}
\]

\subsubsection*{Two-factor square-root models}

A canonical two-factor square-root model is
\[
\begin{aligned}
dY_1(t)
&=\bigl(\mu_1-\lambda_{11}Y_1(t)-\lambda_{12}Y_2(t)\bigr)dt
 +\sqrt{Y_1(t)}\,d\widetilde W_1(t),\\
dY_2(t)
&=\bigl(\mu_2-\lambda_{21}Y_1(t)-\lambda_{22}Y_2(t)\bigr)dt
 +\sqrt{Y_2(t)}\,d\widetilde W_2(t),\\
r_t&=\delta_0+\delta_1Y_1(t)+\delta_2Y_2(t).
\end{aligned}
\]
Conditions such as
\[
\mu_1,\mu_2\geq0,
\qquad
\lambda_{12},\lambda_{21}\leq0
\]
make the drift point inward on the boundary of the nonnegative quadrant and help preserve nonnegativity. Mean reversion is a separate stability property: it requires the drift matrix
\[
\Lambda=
\begin{pmatrix}
\lambda_{11}&\lambda_{12}\\
\lambda_{21}&\lambda_{22}
\end{pmatrix}
\]
to have eigenvalues with positive real parts. Additional Feller-type restrictions are used when strict positivity is required.

A mixed square-root/Gaussian specification is
\[
\begin{aligned}
dY_1(t)
&=\bigl(\mu-\lambda_1Y_1(t)\bigr)dt
 +\sqrt{Y_1(t)}\,d\widetilde W_1(t),\\
dY_2(t)
&=-\lambda_2Y_2(t)dt
 +\sigma_{21}\sqrt{Y_1(t)}\,d\widetilde W_1(t)
 +\sqrt{\alpha+\beta Y_1(t)}\,d\widetilde W_2(t),\\
r_t&=\delta_0+\delta_1Y_1(t)+\delta_2Y_2(t),
\end{aligned}
\]
where $\widetilde W_1$ and $\widetilde W_2$ are independent and the parameters are chosen so that the square roots are well defined.

In affine multi-factor models, the bond price generally has the form
\[
P(t,T)
=\exp\!\left(
 A(t,T)-C_1(t,T)Y_1(t)-C_2(t,T)Y_2(t)
\right),
\]
with deterministic coefficient functions obtained from Riccati equations. The yield is affine in the factors.

\subsection{Heath--Jarrow--Morton models}

The Heath--Jarrow--Morton approach models the full instantaneous-forward curve. Under the physical measure $\mathbb{P}$, suppose
\[
df(t,T)
=\alpha(t,T)dt+\sigma(t,T)^\top dW_t,
\qquad 0\leq t\leq T,
\]
where $W$ is $d$-dimensional and $\sigma(t,T)\in\mathbb{R}^d$.

Define
\[
\alpha^*(t,T)=\int_t^T\alpha(t,u)\,du,
\qquad
\sigma^*(t,T)=\int_t^T\sigma(t,u)\,du.
\]
Let
\[
g(t,T)=\int_t^Tf(t,u)\,du,
\qquad
P(t,T)=e^{-g(t,T)}.
\]
Leibniz's rule gives
\[
dg(t,T)
=\bigl(-r_t+\alpha^*(t,T)\bigr)dt
 +\sigma^*(t,T)^\top dW_t.
\]
It\^o's formula therefore yields
\[
\boxed{
\frac{dP(t,T)}{P(t,T)}
=\left(
 r_t-\alpha^*(t,T)
 +\frac12\|\sigma^*(t,T)\|^2
 \right)dt
-\sigma^*(t,T)^\top dW_t.}
\]

Let $\theta_t\in\mathbb{R}^d$ be the market price of Brownian risk and define
\[
dW_t^{\widetilde{\mathbb{P}}}=dW_t+\theta_tdt.
\]
Under the corresponding equivalent measure, substitution gives
\[
\frac{dP(t,T)}{P(t,T)}
=\left(
 r_t-\alpha^*(t,T)
 +\frac12\|\sigma^*(t,T)\|^2
 +\sigma^*(t,T)^\top\theta_t
 \right)dt
-\sigma^*(t,T)^\top dW_t^{\widetilde{\mathbb{P}}}.
\]
For discounted bond prices to be local martingales, the HJM drift restriction under $\mathbb{P}$ is
\[
\boxed{
\alpha^*(t,T)
=\sigma^*(t,T)^\top\theta_t
 +\frac12\|\sigma^*(t,T)\|^2.}
\]
Differentiating with respect to $T$ gives
\[
\alpha(t,T)
=\sigma(t,T)^\top\theta_t
 +\sigma(t,T)^\top\sigma^*(t,T).
\]
Hence, under the risk-neutral measure,
\[
\boxed{
df(t,T)
=\sigma(t,T)^\top\sigma^*(t,T)dt
 +\sigma(t,T)^\top dW_t^{\widetilde{\mathbb{P}}}.}
\]
The corresponding bond dynamics are
\[
\boxed{
\frac{dP(t,T)}{P(t,T)}
=r_tdt-\sigma^*(t,T)^\top dW_t^{\widetilde{\mathbb{P}}}.}
\]
Thus, once the forward-rate volatility is specified, absence of arbitrage determines the risk-neutral drift.

\section{Jump models}

\subsection{Review of Poisson processes}

Diffusion models have continuous paths and therefore cannot generate sudden price moves. Jump models can capture discontinuities, heavier tails, and some features of implied-volatility smiles and skews.

Let $\xi_1,\xi_2,\ldots$ be independent exponential random variables with parameter $\lambda>0$:
\[
\mathbb{P}(\xi_i>t)=e^{-\lambda t},
\qquad
\mathbb{E}[\xi_i]=\frac1\lambda.
\]
The exponential distribution is memoryless:
\[
\mathbb{P}(\xi_i>s+t\mid\xi_i>s)
=\mathbb{P}(\xi_i>t).
\]
Define the $n$th arrival time by
\[
T_n=\sum_{i=1}^n\xi_i
\]
and the counting process by
\[
N_t=\max\{n\geq0:T_n\leq t\}.
\]
Then $N$ is a Poisson process of intensity $\lambda$ and
\[
\mathbb{P}(N_t=k)
=e^{-\lambda t}\frac{(\lambda t)^k}{k!},
\qquad
\mathbb{E}[N_t]=\operatorname{Var}(N_t)=\lambda t.
\]
The arrival time $T_n$ has the Erlang density
\[
f_{T_n}(s)
=\frac{\lambda^n s^{n-1}}{(n-1)!}e^{-\lambda s},
\qquad s>0.
\]
The compensated process
\[
M_t=N_t-\lambda t
\]
is a martingale.

Let $Y_1,Y_2,\ldots$ be i.i.d., independent of $N$, and suppose $\mathbb{E}|Y_1|<\infty$. The compound Poisson process is
\[
Q_t=\sum_{i=1}^{N_t}Y_i.
\]
A standard compound Poisson process has at most one arrival at any given time almost surely; the jump at an arrival may have a random size $Y_i$. If
\[
\beta=\mathbb{E}[Y_1],
\]
then
\[
Q_t-\lambda\beta t
\]
is a martingale. If $M_Y(u)=\mathbb{E}[e^{uY_1}]$ is finite, then
\[
\mathbb{E}[e^{uQ_t}]
=\exp\!\left(\lambda t\bigl(M_Y(u)-1\bigr)\right).
\]

\subsection{Stochastic calculus for jump processes}

For a compound Poisson process with jump sizes $Y_i$, define the random jump measure by
\[
J((0,t]\times A)
=\sum_{i=1}^{N_t}\mathbf{1}_{\{Y_i\in A\}},
\]
for Borel sets $A\subseteq\mathbb{R}$. Then
\[
Q_t=\int_0^t\int_{\mathbb{R}}y\,J(ds,dy).
\]
If the jump-size distribution is $G(dy)$, the compensator is
\[
\nu(dy)ds=\lambda G(dy)ds,
\]
and the compensated jump measure is
\[
\widetilde J(ds,dy)=J(ds,dy)-\lambda G(dy)ds.
\]
For predictable integrands satisfying the appropriate integrability condition,
\[
\int_0^t\int H_s(y)\,\widetilde J(ds,dy)
\]
is a local martingale.

A finite-activity jump diffusion may be written as
\[
X_t
=X_0+\int_0^tb_sds+\int_0^t\gamma_s^\top dW_s
 +\int_0^t\int_{\mathbb{R}}y\,J(ds,dy).
\]
Its continuous part is
\[
X_t^c
=X_0+\int_0^tb_sds+\int_0^t\gamma_s^\top dW_s,
\]
and its jump at time $s$ is
\[
\Delta X_s=X_s-X_{s-}.
\]
Stochastic integrands are evaluated predictably, using left limits at jumps. For example,
\[
\int_0^tH_s\,dX_s
=\int_0^tH_s b_sds
 +\int_0^tH_s\gamma_s^\top dW_s
 +\sum_{0<s\leq t}H_s\Delta X_s,
\]
where $H$ is predictable. If $X$ is a local martingale and $H$ is $X$-integrable, then $H\mathbin{\boldsymbol\cdot}X$ is a local martingale; predictability alone does not make an integral with respect to an arbitrary semimartingale a martingale.

For semimartingales $X$ and $Y$,
\[
[X,Y]_t
=[X^c,Y^c]_t
 +\sum_{0<s\leq t}\Delta X_s\Delta Y_s.
\]
In particular,
\[
[X]_t=[X^c]_t+\sum_{0<s\leq t}(\Delta X_s)^2.
\]

\begin{thm}[It\^o formula with jumps]
Let $f\in C^2(\mathbb{R})$ and let $X$ be a semimartingale. Then
\[
\begin{aligned}
f(X_t)
&=f(X_0)
 +\int_0^tf'(X_{s-})\,dX_s
 +\frac12\int_0^tf''(X_{s-})\,d[X^c]_s\\
&\quad+\sum_{0<s\leq t}
 \left(
 f(X_s)-f(X_{s-})-f'(X_{s-})\Delta X_s
 \right).
\end{aligned}
\]
\end{thm}

For a $C^{1,2}$ function $f(t,x)$, the time-dependent form is
\[
\begin{aligned}
f(t,X_t)
&=f(0,X_0)
 +\int_0^tf_t(s,X_{s-})ds
 +\int_0^tf_x(s,X_{s-})dX_s\\
&\quad+\frac12\int_0^tf_{xx}(s,X_{s-})d[X^c]_s\\
&\quad+\sum_{0<s\leq t}
 \left(
 f(s,X_s)-f(s,X_{s-})
 -f_x(s,X_{s-})\Delta X_s
 \right).
\end{aligned}
\]

The product rule is
\[
d(X_tY_t)
=X_{t-}\,dY_t+Y_{t-}\,dX_t+d[X,Y]_t.
\]

If $X_0=0$ and $\Delta X_t>-1$, the Dol\'eans--Dade exponential is
\[
\mathcal{E}(X)_t
=\exp\!\left(X_t-\frac12[X^c]_t\right)
 \prod_{0<s\leq t}(1+\Delta X_s)e^{-\Delta X_s}.
\]
It is the unique solution of
\[
dZ_t=Z_{t-}\,dX_t,
\qquad Z_0=1.
\]
When $X$ is a local martingale, $\mathcal{E}(X)$ is a local martingale, but additional integrability is required for it to be a true martingale.

\subsection{Changes of measure for Poisson processes}

Suppose $N$ is Poisson with intensity $\lambda$ under $\mathbb{P}$, and let $\widetilde\lambda>0$. The process
\[
Z_t
=\exp\!\bigl((\lambda-\widetilde\lambda)t\bigr)
 \left(\frac{\widetilde\lambda}{\lambda}\right)^{N_t}
\]
is a martingale. Under the measure defined by
\[
\frac{d\widetilde{\mathbb{P}}}{d\mathbb{P}}
\bigg|_{\mathcal{F}_T}=Z_T,
\]
the process $N$ is Poisson with intensity $\widetilde\lambda$.

Now suppose the marks have density $g$ under $\mathbb{P}$, and choose a target density $\widetilde g$ that is absolutely continuous with respect to $g$; equivalently,
\[
g(y)=0\implies\widetilde g(y)=0
\quad\text{for almost every }y.
\]
Then
\[
Z_t
=\exp\!\bigl((\lambda-\widetilde\lambda)t\bigr)
 \prod_{i=1}^{N_t}
 \frac{\widetilde\lambda\,\widetilde g(Y_i)}
      {\lambda g(Y_i)}
\]
changes the intensity to $\widetilde\lambda$ and the mark density to $\widetilde g$, provided $Z$ is a true martingale. The new measure is equivalent to $\mathbb{P}$ precisely when $g$ and $\widetilde g$ have the same support up to null sets.

A Brownian change of drift can be combined with this jump change of measure. If
\[
Z_t^{W}
=\exp\!\left(
 -\int_0^t\theta_s^\top dW_s
 -\frac12\int_0^t\|\theta_s\|^2ds
\right)
\]
and $Z^J$ is the jump likelihood above, then, under appropriate martingale and dependence assumptions, $Z_t=Z_t^WZ_t^J$ defines a measure under which
\[
\widetilde W_t=W_t+\int_0^t\theta_sds
\]
is Brownian and the jump process has intensity $\widetilde\lambda$ and mark density $\widetilde g$.

\subsection{Jump-diffusion asset models}

Let
\[
Q_t=\sum_{i=1}^{N_t}Y_i,
\qquad Y_i>-1,
\qquad
\beta=\mathbb{E}[Y_i].
\]
Under the physical measure, consider
\[
\frac{dS_t}{S_{t-}}
=\alpha\,dt+\sigma\,dW_t+d(Q_t-\lambda\beta t).
\]
Equivalently,
\[
\frac{dS_t}{S_{t-}}
=(\alpha-\lambda\beta)dt+\sigma dW_t+dQ_t.
\]
Its solution is
\[
S_t
=S_0\exp\!\left(
 (\alpha-\lambda\beta-\tfrac12\sigma^2)t
 +\sigma W_t
\right)
\prod_{i=1}^{N_t}(1+Y_i).
\]

Under an equivalent measure $\widetilde{\mathbb{P}}$, suppose the Brownian kernel is $\theta$, the jump intensity is $\widetilde\lambda$, the mark density is $\widetilde g$, and
\[
\widetilde\beta
=\widetilde{\mathbb{E}}[Y].
\]
With the convention
\[
d\widetilde W_t=dW_t+\theta_tdt,
\]
the discounted stock is a local martingale exactly when
\[
\boxed{
\alpha-\lambda\beta-\sigma\theta_t
=r-\widetilde\lambda\widetilde\beta.}
\]
Equivalently, under $\widetilde{\mathbb{P}}$,
\[
\frac{dS_t}{S_{t-}}
=r\,dt+\sigma\,d\widetilde W_t
 +\int_{(-1,\infty)}y\,\widetilde J(dt,dy),
\]
where
\[
\widetilde J(dt,dy)
=J(dt,dy)-\widetilde\lambda\widetilde g(y)dy\,dt.
\]
In raw-jump form,
\[
\frac{dS_t}{S_{t-}}
=(r-\widetilde\lambda\widetilde\beta)dt
 +\sigma d\widetilde W_t
 +\int_{(-1,\infty)}y\,J(dt,dy).
\]

There are generally many Brownian kernels, jump intensities, and mark densities satisfying the martingale condition. This nonuniqueness is a manifestation of market incompleteness: the stock and money-market account do not usually span all Brownian and jump risks.

\subsection{Pricing derivatives driven by jump processes}

Let $\tau=T-t$. Under $\widetilde{\mathbb{P}}$,
\[
S_T
=S_t\exp\!\left(
(r-\widetilde\lambda\widetilde\beta-\tfrac12\sigma^2)\tau
+\sigma(\widetilde W_T-\widetilde W_t)
\right)
\prod_{i=N_t+1}^{N_T}(1+Y_i).
\]
Define, conditional on $N_T-N_t=j$,
\[
R_j
=e^{-\widetilde\lambda\widetilde\beta\tau}
 \prod_{i=1}^j(1+Y_i),
\qquad R_0=e^{-\widetilde\lambda\widetilde\beta\tau}.
\]
Conditional on the number and sizes of jumps, the remaining uncertainty is Black--Scholes lognormal risk. Therefore the European call price is the Poisson mixture
\[
\boxed{
\begin{aligned}
c(t,S_t)
&=\sum_{j=0}^{\infty}
 e^{-\widetilde\lambda\tau}
 \frac{(\widetilde\lambda\tau)^j}{j!}
 \widetilde{\mathbb{E}}\!\left[
 C_{\mathrm{BS}}(\tau,S_tR_j;K,r,\sigma)
 \right],
\end{aligned}}
\]
where the expectation inside the sum is over $j$ independent marks with density $\widetilde g$. The notation $C_{\mathrm{BS}}$ denotes the Black--Scholes call-pricing function; its arguments are time to maturity $\tau$, spot $s$, strike $K$, rate $r$, and volatility $\sigma$.

\subsubsection*{Partial integro-differential equation}

Assume $c(t,x)$ is sufficiently smooth and has terminal condition
\[
c(T,x)=(x-K)^+.
\]
The risk-neutral generator of the stock in raw-jump form is
\[
\begin{aligned}
\mathcal{L}c(t,x)
&=(r-\widetilde\lambda\widetilde\beta)x c_x(t,x)
 +\frac12\sigma^2x^2c_{xx}(t,x)\\
&\quad+\widetilde\lambda
 \int_{-1}^{\infty}
 \bigl[c(t,x(1+y))-c(t,x)\bigr]
 \widetilde g(y)\,dy.
\end{aligned}
\]
The discounted option value is a martingale, so $c$ solves
\[
\boxed{
\begin{aligned}
0={}&c_t(t,x)
 +(r-\widetilde\lambda\widetilde\beta)x c_x(t,x)
 +\frac12\sigma^2x^2c_{xx}(t,x)-rc(t,x)\\
&+\widetilde\lambda
 \int_{-1}^{\infty}
 \bigl[c(t,x(1+y))-c(t,x)\bigr]
 \widetilde g(y)\,dy,
\end{aligned}}
\]
with $c(T,x)=(x-K)^+$.

Equivalently, using the compensated jump measure, the PIDE may be written
\[
\begin{aligned}
0={}&c_t+rx c_x+\frac12\sigma^2x^2c_{xx}-rc\\
&+\widetilde\lambda
 \int_{-1}^{\infty}
 \bigl[c(t,x(1+y))-c(t,x)-xyc_x(t,x)\bigr]
 \widetilde g(y)\,dy.
\end{aligned}
\]

\subsubsection*{Delta hedging and incompleteness}

Suppose a self-financing portfolio holds $\Gamma_t$ shares of stock. In discounted terms,
\[
\begin{aligned}
d(e^{-rt}X_t)
&=e^{-rt}\Gamma_t\sigma S_{t-}\,d\widetilde W_t\\
&\quad+e^{-rt}\Gamma_tS_{t-}
 \int_{-1}^{\infty}y\,\widetilde J(dt,dy).
\end{aligned}
\]
The discounted option value satisfies
\[
\begin{aligned}
d(e^{-rt}c(t,S_t))
&=e^{-rt}\sigma S_{t-}c_x(t,S_{t-})\,d\widetilde W_t\\
&\quad+e^{-rt}
 \int_{-1}^{\infty}
 \bigl[c(t,S_{t-}(1+y))-c(t,S_{t-})\bigr]
 \widetilde J(dt,dy).
\end{aligned}
\]
Choosing
\[
\Gamma_t=c_x(t,S_{t-})
\]
removes the Brownian mismatch, but the jump mismatch remains:
\[
\begin{aligned}
d\bigl(e^{-rt}(X_t-c(t,S_t))\bigr)
=e^{-rt}\int_{-1}^{\infty}H_t(y)\,\widetilde J(dt,dy),
\end{aligned}
\]
where
\[
H_t(y)
=S_{t-}y\,c_x(t,S_{t-})
-\bigl[c(t,S_{t-}(1+y))-c(t,S_{t-})\bigr].
\]
For a convex option value,
\[
c(t,x(1+y))-c(t,x)
\geq xy\,c_x(t,x),
\]
so
\[
H_t(y)\leq0.
\]
Consequently, between jumps the compensator makes the delta-hedging error drift upward, so the stock hedge tends to outperform the option. At a jump, the error changes by $H_t(Y)\leq0$, so the hedge underperforms. A single stock position cannot in general eliminate both Brownian and jump-size risk; this is another direct demonstration that the jump-diffusion market is incomplete.
